\documentclass[../../../include/open-logic-section]{subfiles}

\begin{document}

\olfileid[fr]{sth}{ord-arithmetic}{expo}
\olsection{Exponentiation ordinale}

Passons maintenant à l'exponentiation ordinale.
Malheureusement, il n'en existe pas de définition
synthétique \emph{simple}.

Il existe certes des définitions synthétiques
\emph{explicites}. En voici une. Soit
$\text{finfun}(\alpha,\beta)$ l'ensemble des
fonctions $f \colon \beta \to \alpha$
pour lesquelles
$\Setabs{\gamma \in \beta}{f(\gamma) \neq 0}$
est équipotent à un nombre naturel.\footnote{Note
de l'édition française : l'original inverse ici
le domaine et le codomaine. Pour représenter
$\ordexpo{\alpha}{\beta}$, il faut des fonctions
de $\beta$ vers $\alpha$, dont le support est
un sous-ensemble fini de $\beta$.}
Définissons un bon ordre sur
$\text{finfun}(\alpha,\beta)$ en posant
$f \sqsubset g$ si et seulement si
$f \neq g$ et
$f(\gamma_0) < g(\gamma_0)$, où
$\gamma_0 =
\text{max}\Setabs{\gamma \in \beta}
{f(\gamma) \neq g(\gamma)}$.
Nous pouvons alors définir
$\ordexpo{\alpha}{\beta}$ comme
$\ordtype{\text{finfun}(\alpha, \beta),
\sqsubset}$. Potter emploie cette définition
explicite, puis explique aussitôt :
\begin{quote}
	Le choix de cet ordre est entièrement déterminé
	par notre souhait d'obtenir une définition de
	l'exponentiation ordinale qui satisfasse la
	condition de récursion appropriée\ldots ;
	et il est bien plus difficile à se représenter
	que la somme ordonnée ou le produit ordonné.
	\citep[p.~199]{Potter2004}
\end{quote}
En effet. Nous avons expliqué l'addition comme
«~une copie de $\alpha$ suivie d'une copie de
$\beta$~», et la multiplication comme
«~une suite de $\beta$ copies de $\alpha$~».
Mais nous n'avons rien d'aussi concis à dire sur
$\text{finfun}(\alpha, \gamma)$. Nous définirons
donc l'exponentiation ordinale \emph{par}
récursion transfinie :

\begin{defn}\ollabel{ordexporecursion}
Pour $\alpha>0$, posons :
\begin{align*}
	\ordexpo{\alpha}{0} &= 1\\
	\ordexpo{\alpha}{\beta\ordplus 1}
	&= \ordexpo{\alpha}{\beta} \ordtimes \alpha\\
	\ordexpo{\alpha}{\beta}
	&= \bigcup_{\delta < \beta}
	\ordexpo{\alpha}{\delta}
	&& \text{si $\beta$ est un ordinal limite}
\end{align*}
Pour la base nulle, posons
$\ordexpo{0}{0}=1$ et
$\ordexpo{0}{\beta}=0$ si $\beta>0$.
\footnote{Note de l'édition française : la règle
pour les exposants limites, appliquée dans
l'original sans restriction à la base $0$,
donnerait $\ordexpo{0}{\omega}=1$ à cause du
terme d'exposant $0$. La restriction et le cas
de base nulle rétablissent la convention
ordinale usuelle.}
\end{defn}

Si nous travaillions \emph{en tant que} théoriciens
des ensembles, nous pourrions étudier diverses
propriétés de l'exponentiation ordinale. Nous
n'ajouterons guère qu'un fait peu surprenant :
elle n'est pas commutative. Ainsi,
$\ordexpo{2}{\omega} =
\bigcup_{\delta < \omega}\ordexpo{2}{\delta}
= \omega$, tandis que
$\ordexpo{\omega}{2} = \omega \ordtimes \omega$.
Mais nous ne devrions \emph{pas nous attendre}
à ce que l'exponentiation soit commutative,
puisqu'elle ne l'est déjà pas pour les nombres
naturels :
$\ordexpo{2}{3} = 8 < 9 = \ordexpo{3}{2}$.

\begin{prob}
À l'aide de la récurrence transfinie, démontrez
que, si nous définissons
$\ordexpo{\alpha}{\beta} =
\ordtype{\text{finfun}(\alpha, \beta),
\sqsubset}$, nous obtenons les équations
de récursion de la
\olref[sth][ord-arithmetic][expo]{ordexporecursion}.
\end{prob}

\end{document}
